% !TEX program = pdflatex
% TD1 - Suites et limites de suites (Terminale, maths complémentaires) -- SQUELETTE
% Compilation : pdflatex TD-suites-limites.tex
\documentclass[11pt,a4paper]{article}

\usepackage[utf8]{inputenc}
\usepackage[T1]{fontenc}
\usepackage[french]{babel}
\usepackage{lmodern}
\usepackage{amsmath,amssymb,amsfonts}
\usepackage[a4paper,margin=2.2cm]{geometry}
\usepackage{enumitem}
\usepackage{fancyhdr}
\usepackage{xcolor}
\usepackage{array}
\usepackage{graphicx}       % pour insérer les captures
% \graphicspath{{images/}}  % décommenter si les captures sont dans un sous-dossier

%------------------------------------------------------------
% Mise en forme
%------------------------------------------------------------
\definecolor{bleu}{RGB}{30,80,140}

\pagestyle{fancy}
\fancyhf{}
\lhead{\small Mathématiques -- Terminale}
\rhead{\small TD1 : Suites et limites de suites}
\cfoot{\small \thepage}
\renewcommand{\headrulewidth}{0.4pt}

% Environnement "exercice" numéroté avec titre
\newcounter{exo}
\newcommand{\exo}[1]{%
  \refstepcounter{exo}%
  \vspace{0.8cm}
  \noindent{\color{bleu}\large\bfseries Exercice \theexo\ -- #1}\par
  \nopagebreak\smallskip
}

% Raccourci pour insérer une capture (fichier + largeur en fraction de page)
% Utilisation :  \capture{nom-fichier.png}{0.6}
\newcommand{\capture}[2]{%
  \begin{center}\includegraphics[width=#2\textwidth]{#1}\end{center}}

\setlength{\parindent}{0pt}
\setlist[enumerate]{leftmargin=*,itemsep=4pt,topsep=4pt}

%------------------------------------------------------------
\begin{document}

%------------------- En-tête -------------------
\begin{center}
  {\color{bleu}\LARGE\bfseries TD1 -- Suites et limites de suites}\\[4pt]
  {\large Terminale -- Mathématiques complémentaires}
\end{center}

\vspace{6pt}
\noindent\fbox{\parbox{\dimexpr\textwidth-2\fboxsep-2\fboxrule}{%
  \small \textbf{Objectifs du TD :} s'entraîner sur les calculs de limites avec différents
  outils.
  
  \vspace{5pt}
  \textbf{Ces exercices sont extraits du manuel en ligne sésamaths : page 31.}
  
  \vspace{5pt}
  \textbf{\underline{Exercices à faire sur une feuille ou un cahier séparé.}}}}

%============================================================
%  EXERCICE 1  -- révisions de première
%============================================================
\exo{Représentation graphique et calculs de limites}
% --- Énoncé ---
\begin{minipage}[t]{.5\textwidth}%
    \vspace{0pt}
    \capture{exo1TD1.png}{1}   % <-- insérer une capture ici
\end{minipage}%
\hfill
\begin{minipage}[t]{.5\textwidth}%
    \vspace{0pt}
    \capture{exo2TD1.png}{1}   % <-- insérer une capture ici

    \vspace{10pt}
    \capture{exo3TD1.png}{1}   % <-- insérer une capture ici

\end{minipage}%


%============================================================
%  EXERCICE 2  -- conjecturer une limite (tableau / graphique)
%============================================================
\exo{Calculs de limites avec les théorèmes}
% --- Énoncé ---
\begin{minipage}[t]{.5\textwidth}%
    \vspace{0pt}
    \capture{exo4TD1.png}{1}   % <-- insérer une capture ici
\end{minipage}%
\hfill
\begin{minipage}[t]{.5\textwidth}%
    \vspace{0pt}
    \capture{exo5TD1.png}{1}   % <-- insérer une capture ici
\end{minipage}%


%============================================================
%  EXERCICE 2  -- conjecturer une limite (tableau / graphique)
%============================================================
\exo{Une suite pas comme les autres}
% --- Énoncé ---
\begin{minipage}[t]{.5\textwidth}%
    \vspace{0pt}
    \capture{exo6.png}{1}   % <-- insérer une capture ici
\end{minipage}%
\hfill
\begin{minipage}[t]{.5\textwidth}%
    \vspace{0pt}
    % \capture{exo5TD1.png}{1}   % <-- insérer une capture ici
\end{minipage}%


%============================================================
%  (dupliquer un bloc \exo{...} pour ajouter un exercice)
%============================================================

\end{document}
